% Generated draft: requires mathematical and source review.
\documentclass[11pt]{article}
\usepackage[margin=25mm]{geometry}
\usepackage{fontspec}
\setmainfont{Latin Modern Roman}
\usepackage{amsmath,amssymb,mathtools,unicode-math}
\setmathfont{Latin Modern Math}
\usepackage{enumitem,needspace,xurl,hyperref}
\setlength{\parindent}{0pt}
\setlength{\parskip}{5pt}
\setlength{\emergencystretch}{3em}
\newcommand{\N}{\mathbb{N}}
\newcommand{\Z}{\mathbb{Z}}
\newcommand{\Q}{\mathbb{Q}}
\newcommand{\R}{\mathbb{R}}
\newcommand{\C}{\mathbb{C}}
\newcommand{\F}{\mathbb{F}}
\newcommand{\A}{\mathbb{A}}
\newcommand{\PP}{\mathbb{P}}
\newcommand{\E}{\mathbb{E}}
\newcommand{\eps}{\varepsilon}
\newcommand{\dd}{\,\mathrm{d}}
\begin{document}
\section*{UnsolvedMath ID: 6200085}
\textbf{Source problem code:} catalogue entry
\subsection{Definitions and mathematical statement}
Take the finite presentation appearing in the source,
\[
 G=\langle g_1,\ldots,g_k\mid R_1,\ldots,R_m\rangle,
 \qquad \ell_j=|R_j|,
\]
where word length counts letters in the generators and their inverses. For a representation \(\rho:G\to\operatorname{Isom}(\mathbb H^n)\), with the usual curvature-\(-1\) hyperbolic metric \(d\), define
\[
 B_n([\rho])=\inf_{x\in\mathbb H^n}\max_{1\le i\le k}
 d(x,\rho(g_i)x).
\]
Conjugating \(\rho\) by a hyperbolic isometry leaves this value unchanged, explaining the notation \([\rho]\). No attainment of the infimum is assumed.

The source asks for a constructive proof of an ``above theorem'', more precisely an explicit finite constant depending only on the displayed presentation parameters,
\[
 C=C(n,k,m,\ell_1,\ldots,\ell_m),\qquad
 \forall[\rho]\in D_n(G):\quad B_n([\rho])\le C.
\]
The bound is independent of the individual representation and uses the chosen generating set. ``Explicit'' should in particular permit evaluation from the allowed finite data, rather than merely asserting existence of an unspecified supremum.

The referenced theorem is absent, as are its hypotheses on \(G\) and the definition of \(D_n(G)\). Thus the admissible representations and group class cannot be inferred: \(D_n(G)\) is not silently replaced by all representations or even by all discrete faithful representations. The missing restrictions are essential to a boundedness claim. The word-length convention and range of dimensions also need confirmation in the original source.

The summary labels the problem open and mentions domains of discontinuity, stable translation, Margulis invariants, and effective linearization, without giving a precise connection to the displayed bound or a theorem citation. These descriptions do not restore the missing theorem. The historical problem label and later open assessment remain unverified source information.
\subsection{Short English statement}
Give a computable uniform bound, in terms of dimension and presentation sizes, for the least possible maximal displacement of the generators over the intended deformation space. The preceding theorem and the definition of that deformation space are missing, so the admissible groups and representations cannot yet be fixed.
\subsection{Sources}
\begin{itemize}
\item Local registry: \url{data/unsolvedmath/problems.json}, permanent UnsolvedMath ID \(6200085\), source code catalogue entry; supplied statement, background, and research summary, updated 2026-08-06.
\item Kapovich, \emph{Problems on Boundaries of Groups and Kleinian Groups} (2005), Problem 85, PDF page 21: \url{https://www.math.ucdavis.edu/~kapovich/EPR/problems.pdf}. Supplied normalized PDF-text extraction accessed 2026-07-29; current status marked as needing review.
\item The ``above theorem'' and the definition of \(D_n(G)\) are referenced but not included. The summary supplies no separate constructive-bound or effective-linearization publication.
\end{itemize}
\end{document}
